\documentclass[11pt]{article}
\usepackage[T1]{fontenc}
\usepackage{lmodern}
\usepackage{amsmath,amssymb,amsthm,mathtools}
\usepackage{microtype}
\usepackage[a4paper,margin=30mm]{geometry}
\usepackage{booktabs,longtable,array}
\usepackage{enumitem}
\usepackage{xcolor}
\usepackage{hyperref}
\hypersetup{colorlinks=true,linkcolor=blue!45!black,citecolor=green!35!black,urlcolor=blue!55!black,pdfauthor={OpenAI},pdftitle={q-Pochhammer Symbols and q-Binomial Coefficients}}
\setcounter{tocdepth}{2}
\numberwithin{equation}{section}
\newtheorem{theorem}{Theorem}[section]
\newtheorem{lemma}[theorem]{Lemma}
\newtheorem{proposition}[theorem]{Proposition}
\newtheorem{corollary}[theorem]{Corollary}
\newtheorem{identity}[theorem]{Identity}
\theoremstyle{definition}
\newtheorem{definition}[theorem]{Definition}
\newtheorem{example}[theorem]{Example}
\newtheorem{remark}[theorem]{Remark}
\newcommand{\qbinom}[2]{\genfrac{[}{]}{0pt}{}{#1}{#2}_{q}}
\newcommand{\qmult}[1]{\genfrac{[}{]}{0pt}{}{#1}_{q}}
\newcommand{\N}{\mathbb N}
\newcommand{\Z}{\mathbb Z}
\newcommand{\C}{\mathbb C}
\newcommand{\F}{\mathbb F}
\newcommand{\Li}{\operatorname{Li}}
\newcommand{\CT}{\operatorname{CT}}
\newcommand{\ord}{\operatorname{ord}}
\newcommand{\inv}{\operatorname{inv}}
\newcommand{\area}{\operatorname{area}}
\newcommand{\eps}{\varepsilon}
\title{\bfseries q-Pochhammer Symbols and q-Binomial Coefficients\\[3mm]
Algebra, Combinatorics, Analysis, Hypergeometric Summation,\\Finite Geometry, Roots of Unity, and q-Calculus}
\author{A self-contained mathematical survey}
\date{August 2026}
\begin{document}
\maketitle
\begin{abstract}
The finite and infinite $q$-Pochhammer symbols and the Gaussian, or $q$-binomial, coefficients form a compact language in which a remarkable part of algebraic combinatorics, partition theory, finite geometry, $q$-calculus, and basic hypergeometric analysis can be expressed.  This article develops that language from first principles.  It proves the elementary product laws, the finite and infinite $q$-binomial theorems, Euler's expansions, $q$-binomial inversion, convolution identities, Heine's transformation and the $q$-Gauss and terminating $q$-Chu--Vandermonde summations.  It gives mutually translating models by inversions of words, areas of lattice paths, partitions in rectangles, subspaces over finite fields, flags, symmetric functions, and weighted selections.  It also treats cyclotomic factorization, the $q$-Lucas theorem, cyclic sieving for subsets, asymptotics near $q=1$, the $q$-gamma function, Jackson differentiation and integration, and the quantum-plane binomial theorem.  Except where a result is explicitly identified as contextual, proofs are supplied in full.
\end{abstract}
\tableofcontents

\section{Conventions and three complementary viewpoints}
Throughout, $q$ is either an indeterminate or a complex number.  Algebraic identities are first interpreted in the rational-function field $\C(q)$, or in a formal power-series ring such as $\C(q)[[z]]$.  Analytic assertions involving infinite products are made under $|q|<1$ unless another domain is stated.  Empty products equal $1$, and a Gaussian coefficient is understood to be zero when its lower index is outside $0,\ldots,n$.

Three viewpoints will recur.  The product viewpoint sees $q$-Pochhammer symbols as shifted products.  The enumerative viewpoint sees a Gaussian coefficient as a generating polynomial for a statistic.  The analytic viewpoint regards ratios of infinite products as solutions of first-order $q$-difference equations.  Most important identities become nearly inevitable after moving to the right viewpoint.

\section{Finite q-shifted factorials}
\begin{definition}[Finite $q$-Pochhammer symbol]
For $n\in\N$,
\[
 (a;q)_n:=\prod_{j=0}^{n-1}(1-aq^j),\qquad (a;q)_0:=1.
\]
For several parameters we abbreviate
\[
 (a_1,\ldots,a_r;q)_n:=\prod_{i=1}^r(a_i;q)_n.
\]
The same object is called a finite $q$-shifted factorial.
\end{definition}

\begin{proposition}[Concatenation, quotient, and splitting]
For nonnegative integers $m,n$,
\begin{align}
(a;q)_{m+n}&=(a;q)_m(aq^m;q)_n,\label{eq:concat}\\
\frac{(a;q)_m}{(a;q)_n}&=(aq^n;q)_{m-n}\quad(m\ge n),\label{eq:quot}\\
(a;q)_{mn}&=\prod_{r=0}^{m-1}(aq^r;q^m)_n.\label{eq:split}
\end{align}
More generally, if $N=mn+s$ with $0\le s<m$, then
\[
(a;q)_N=\prod_{r=0}^{s-1}(aq^r;q^m)_{n+1}
       \prod_{r=s}^{m-1}(aq^r;q^m)_n.
\]
\end{proposition}
\begin{proof}
Equation \eqref{eq:concat} partitions the exponent set $\{0,\ldots,m+n-1\}$ into $\{0,\ldots,m-1\}$ and $\{m,\ldots,m+n-1\}$.  Division gives \eqref{eq:quot}.  For \eqref{eq:split}, write each exponent $0\le j<mn$ uniquely as $j=r+m\ell$ with $0\le r<m$ and $0\le\ell<n$.  The final formula is the same residue-class decomposition when the residue classes $0,\ldots,s-1$ contain one additional exponent.
\end{proof}

\begin{proposition}[Reversal and inversion of the base]
For $n\ge0$ and nonzero $a,q$,
\begin{align}
(a;q)_n&=(-a)^nq^{\binom n2}(q^{1-n}/a;q)_n,\label{eq:reverse}\\
(a;q^{-1})_n&=(-a)^nq^{-\binom n2}(a^{-1};q)_n.\label{eq:baseinv}
\end{align}
\end{proposition}
\begin{proof}
Factor $-aq^j$ from $1-aq^j$ and reverse the remaining exponents:
\[
\prod_{j=0}^{n-1}(1-aq^j)
=(-a)^nq^{0+\cdots+(n-1)}
 \prod_{j=0}^{n-1}(1-a^{-1}q^{-j}).
\]
The last exponent set is $\{0,-1,\ldots,1-n\}$ and hence gives \eqref{eq:reverse}.  Replacing $q$ by $q^{-1}$ in the defining product and factoring $-aq^{-j}$ gives \eqref{eq:baseinv}.
\end{proof}

\begin{definition}[Integral index]
Whenever no denominator vanishes, define
\[
(a;q)_{-n}:=\frac1{(aq^{-n};q)_n}\qquad(n\ge1).
\]
Then $(a;q)_{m+n}=(a;q)_m(aq^m;q)_n$ holds for all integral $m,n$ for which both sides are defined.
\end{definition}

\begin{proposition}[Explicit negative-index formula]
For $n\ge0$,
\[
(a;q)_{-n}=\frac{(-q/a)^nq^{\binom n2}}{(q/a;q)_n}.
\]
\end{proposition}
\begin{proof}
Factoring $-aq^{-n+j}$ from the $j$th factor of $(aq^{-n};q)_n$ gives
\[
(aq^{-n};q)_n=(-a)^nq^{-n(n+1)/2}(q/a;q)_n.
\]
Taking reciprocals proves the formula.
\end{proof}

\begin{proposition}[Roots and degree]
As a polynomial in $a$, $(a;q)_n$ has degree $n$, constant term $1$, leading coefficient $(-1)^nq^{\binom n2}$, and roots $q^{-j}$ for $0\le j<n$.
\end{proposition}
\begin{proof}
Every assertion follows directly by inspecting the $n$ linear factors.
\end{proof}

\section{Infinite products and analytic control}
\begin{definition}
For $|q|<1$,
\[
(a;q)_\infty:=\prod_{j=0}^{\infty}(1-aq^j),\qquad
(a_1,\ldots,a_r;q)_\infty:=\prod_i(a_i;q)_\infty.
\]
\end{definition}

\begin{theorem}[Convergence and zeros]
For fixed $|q|<1$, $(a;q)_\infty$ converges locally uniformly as a function of $a$ and defines an entire function whose zeros are exactly $a=q^{-j}$, $j\ge0$, all simple when $q\ne0$.
\end{theorem}
\begin{proof}
On $|a|\le R$, the series $\sum_{j\ge0}|aq^j|$ is bounded by $R/(1-|q|)$.  The standard convergence criterion for products $\prod(1+u_j)$, applied uniformly on compact sets after omitting finitely many factors if necessary, yields local uniform convergence.  Local uniform limits of polynomials are holomorphic, so the result is entire.  If $a=q^{-m}$ one factor vanishes.  Conversely, if no factor vanishes, the tail product is nonzero because $\sum|aq^j|<\infty$, and the finite initial product is nonzero.  At $a=q^{-m}$ only the $m$th factor vanishes and its derivative is $-q^m\ne0$.
\end{proof}

\begin{proposition}[Tail quotient and $q$-difference equation]
For $n\ge0$,
\[
(a;q)_n=\frac{(a;q)_\infty}{(aq^n;q)_\infty},\qquad
(a;q)_\infty=(1-a)(aq;q)_\infty.
\]
\end{proposition}
\begin{proof}
Cancel the common tail of the two convergent products.  The second formula is the case $n=1$.
\end{proof}

\begin{theorem}[Logarithmic expansion]
If $|q|<1$ and $|a|<1$, then
\[
\log(a;q)_\infty=-\sum_{r=1}^{\infty}\frac{a^r}{r(1-q^r)}.
\]
The series converges absolutely and locally uniformly in $|a|<1$.
\end{theorem}
\begin{proof}
Use $\log(1-w)=-\sum_{r\ge1}w^r/r$ in every factor.  Absolute convergence of the double series follows from
\[
\sum_{j\ge0}\sum_{r\ge1}\frac{|a|^r|q|^{jr}}r
=\sum_{r\ge1}\frac{|a|^r}{r(1-|q|^r)}<\infty.
\]
Fubini's theorem then permits interchange of the sums, and $\sum_{j\ge0}q^{jr}=1/(1-q^r)$.
\end{proof}

\section{q-integers, q-factorials, and Gaussian coefficients}
\begin{definition}
For $n\ge0$,
\[
[n]_q:=\frac{1-q^n}{1-q}=1+q+\cdots+q^{n-1},\qquad
[n]_q!:=\prod_{j=1}^n[j]_q,
\]
with $[0]_q!=1$.  For $0\le k\le n$,
\[
\qbinom nk:=\frac{(q;q)_n}{(q;q)_k(q;q)_{n-k}}
=\frac{[n]_q!}{[k]_q![n-k]_q!}.
\]
\end{definition}

\begin{theorem}[The two $q$-Pascal recurrences]
For all integers $n,k$,
\begin{align}
\qbinom nk&=\qbinom{n-1}k+q^{n-k}\qbinom{n-1}{k-1},\label{eq:pascal1}\\
\qbinom nk&=q^k\qbinom{n-1}k+\qbinom{n-1}{k-1}.\label{eq:pascal2}
\end{align}
\end{theorem}
\begin{proof}
For $1\le k\le n-1$, put the two terms in \eqref{eq:pascal1} over the common denominator $(q;q)_k(q;q)_{n-k}$:
\[
\frac{(q;q)_{n-1}\bigl((1-q^{n-k})+q^{n-k}(1-q^k)\bigr)}{(q;q)_k(q;q)_{n-k}}
=\frac{(q;q)_{n-1}(1-q^n)}{(q;q)_k(q;q)_{n-k}}.
\]
This is $\qbinom nk$.  The boundary cases are immediate.  The proof of \eqref{eq:pascal2} is identical, using $q^k(1-q^{n-k})+(1-q^k)=1-q^n$.
\end{proof}

\begin{corollary}[Polynomiality and positivity]
The rational expression defining $\qbinom nk$ is in fact a polynomial in $q$ with nonnegative integral coefficients.
\end{corollary}
\begin{proof}
Induct on $n$ using either $q$-Pascal recurrence, beginning with $\qbinom00=1$.  Addition and multiplication by a monomial preserve nonnegative integral coefficients.
\end{proof}

\begin{proposition}[Symmetry, reciprocity, and degree]
For $0\le k\le n$,
\begin{align}
\qbinom nk&=\qbinom n{n-k},\label{eq:sym}\\
\genfrac{[}{]}{0pt}{}{n}{k}_{q^{-1}}&=q^{-k(n-k)}\qbinom nk.\label{eq:reciprocity}
\end{align}
Consequently $\qbinom nk$ is monic of degree $k(n-k)$ and its coefficients are palindromic.
\end{proposition}
\begin{proof}
Symmetry is immediate from the factorial quotient.  Apply \eqref{eq:baseinv} to each factor $(q^{-1};q^{-1})_r$ in the quotient to obtain \eqref{eq:reciprocity}; the powers simplify to $-k(n-k)$.  Polynomiality and the constant term $1$ then show that reciprocity reverses the coefficient sequence, giving degree $k(n-k)$ and leading coefficient $1$.
\end{proof}

\section{Cyclotomic factorization and a carry criterion}
Let $\Phi_d(q)$ be the $d$th cyclotomic polynomial.  Since $1-q^m=\prod_{d\mid m}\Phi_d(q)$ up to the harmless sign convention at $d=1$, counting multiples gives the following exact factorization.

\begin{theorem}[Cyclotomic factorization]
For $0\le k\le n$,
\[
\qbinom nk=\prod_{d=2}^{n}\Phi_d(q)^{\epsilon_d(n,k)},
\qquad
\epsilon_d(n,k)=\left\lfloor\frac nd\right\rfloor-
\left\lfloor\frac kd\right\rfloor-
\left\lfloor\frac{n-k}d\right\rfloor.
\]
Moreover $\epsilon_d(n,k)\in\{0,1\}$; it equals $1$ exactly when
\[
(k\bmod d)+((n-k)\bmod d)\ge d,
\]
equivalently when addition of the two least base-$d$ digits produces a carry.
\end{theorem}
\begin{proof}
The exponent of $\Phi_d$ in $(q;q)_r=\prod_{j=1}^r(1-q^j)$ is the number of multiples of $d$ among $1,\ldots,r$, namely $\lfloor r/d\rfloor$.  Subtracting denominator exponents gives the displayed formula.  Write $k=ad+r$ and $n-k=bd+s$ with $0\le r,s<d$.  Then $n=(a+b)d+r+s$, so
\[
\epsilon_d=\left\lfloor\frac{r+s}{d}\right\rfloor,
\]
which is $0$ or $1$ and has the asserted carry interpretation.
\end{proof}

\begin{corollary}
Every Gaussian coefficient is a product of distinct cyclotomic polynomials.  In particular it has no repeated complex zero.
\end{corollary}
\begin{proof}
Every exponent in the factorization is $0$ or $1$, and distinct cyclotomic polynomials are coprime in characteristic zero.
\end{proof}

\section{The finite q-binomial theorem}
\begin{theorem}[Finite $q$-binomial theorem]
For $n\ge0$,
\begin{equation}
\prod_{j=0}^{n-1}(1+zq^j)
=\sum_{k=0}^{n}q^{\binom k2}\qbinom nk z^k.\label{eq:fqbt}
\end{equation}
Equivalently,
\[
(z;q)_n=\sum_{k=0}^n(-1)^kq^{\binom k2}\qbinom nk z^k.
\]
\end{theorem}
\begin{proof}
Induct on $n$.  The claim is clear at $n=0$.  Multiplying the $n$th identity by $1+zq^n$, the coefficient of $z^k$ becomes
\[
q^{\binom k2}\qbinom nk+q^nq^{\binom{k-1}2}\qbinom n{k-1}
=q^{\binom k2}\left(\qbinom nk+q^{n+1-k}\qbinom n{k-1}\right),
\]
which equals $q^{\binom k2}\qbinom{n+1}k$ by \eqref{eq:pascal1}.
\end{proof}

\begin{corollary}[Elementary and complete symmetric specializations]
For $m,r\ge0$,
\begin{align}
e_r(1,q,\ldots,q^m)&=q^{\binom r2}\genfrac{[}{]}{0pt}{}{m+1}{r}_q,\label{eq:elementary}\\
h_r(1,q,\ldots,q^m)&=\genfrac{[}{]}{0pt}{}{m+r}{r}_q.\label{eq:complete}
\end{align}
\end{corollary}
\begin{proof}
The generating function for elementary symmetric polynomials is $\prod_{j=0}^m(1+tq^j)$, so \eqref{eq:elementary} is coefficient extraction in \eqref{eq:fqbt}.  For complete symmetric polynomials,
\[
\sum_{r\ge0}h_r(1,q,\ldots,q^m)t^r=\prod_{j=0}^m\frac1{1-tq^j}.
\]
The coefficient formula \eqref{eq:complete} follows by induction on $m+r$ from multiplication by $(1-tq^m)^{-1}$, or from Euler's reciprocal expansion proved below with a finite tail quotient.
\end{proof}

\section{Words, inversions, paths, and partitions}
\begin{theorem}[Binary-word model]
Let $W_{n,k}$ be the set of binary words of length $n$ with $k$ ones, and let
\[
\inv(w)=\#\{(i,j):i<j,\ w_i=1,\ w_j=0\}.
\]
Then
\[
\qbinom nk=\sum_{w\in W_{n,k}}q^{\inv(w)}.
\]
\end{theorem}
\begin{proof}
Let the right side be $F_{n,k}$.  If the last letter is $1$, deleting it gives a word in $W_{n-1,k-1}$ and creates no inversion.  If the last letter is $0$, deleting it gives a word in $W_{n-1,k}$ and the deleted zero forms an inversion with each of the $k$ ones, contributing $q^k$.  Hence
\[
F_{n,k}=F_{n-1,k-1}+q^kF_{n-1,k},
\]
with the same boundary values as \eqref{eq:pascal2}.  Induction proves equality.
\end{proof}

\begin{theorem}[Partitions in a rectangle]
For nonnegative $r,s$,
\[
\genfrac{[}{]}{0pt}{}{r+s}{r}_q
=\sum_{\lambda\subseteq s^r}q^{|\lambda|},
\]
where $\lambda\subseteq s^r$ means that $\lambda$ has at most $r$ parts, each at most $s$.
\end{theorem}
\begin{proof}
Trace the southeast boundary of the Ferrers diagram of $\lambda$ inside an $r$ by $s$ rectangle, encoding a vertical step by $1$ and a horizontal step by $0$.  This is a bijection to binary words with $r$ ones and $s$ zeros.  A cell of the Ferrers diagram corresponds exactly to a vertical step occurring before a horizontal step, so $|\lambda|=\inv(w)$.  Apply the preceding theorem.
\end{proof}

\begin{corollary}[Palindromicity by complementation]
The coefficient of $q^j$ in $\genfrac{[}{]}{0pt}{}{r+s}{r}_q$ equals the coefficient of $q^{rs-j}$.
\end{corollary}
\begin{proof}
Complement a Ferrers diagram in the $r$ by $s$ rectangle and rotate it by $180^\circ$.  This involution replaces its size by $rs-|\lambda|$.
\end{proof}

\section{Finite vector spaces and flags}
In this section $Q$ denotes a prime power, to distinguish it from a formal variable.

\begin{theorem}[Subspace-counting theorem]
The number of $k$-dimensional subspaces of $\F_Q^n$ is $\genfrac{[}{]}{0pt}{}{n}{k}_{Q}$.
\end{theorem}
\begin{proof}
Count ordered linearly independent $k$-tuples in $\F_Q^n$.  There are
\[
(Q^n-1)(Q^n-Q)\cdots(Q^n-Q^{k-1})
\]
such tuples.  Every $k$-dimensional subspace has
\[
(Q^k-1)(Q^k-Q)\cdots(Q^k-Q^{k-1})
\]
ordered bases.  Division, followed by cancellation of $Q^{0+\cdots+(k-1)}$, gives
\[
\prod_{j=0}^{k-1}\frac{Q^{n-j}-1}{Q^{k-j}-1}
=\frac{(Q;Q)_n}{(Q;Q)_k(Q;Q)_{n-k}}.
\]
\end{proof}

\begin{definition}[Gaussian multinomial]
For $k_1+\cdots+k_r=n$,
\[
\genfrac{[}{]}{0pt}{}{n}{k_1,\ldots,k_r}_q
:=\frac{(q;q)_n}{(q;q)_{k_1}\cdots(q;q)_{k_r}}.
\]
\end{definition}

\begin{proposition}[Factorization and flags]
For partial sums $K_j=k_1+\cdots+k_j$,
\[
\genfrac{[}{]}{0pt}{}{n}{k_1,\ldots,k_r}_q
=\prod_{j=1}^{r-1}\genfrac{[}{]}{0pt}{}{n-K_{j-1}}{k_j}_q.
\]
At $q=Q$, this counts flags
\[
0=V_0\subset V_1\subset\cdots\subset V_r=\F_Q^n,
\qquad \dim(V_j/V_{j-1})=k_j.
\]
\end{proposition}
\begin{proof}
The product identity follows by telescopic cancellation of the factorials.  To count flags, first choose $V_1$, then choose $V_2/V_1$ as a $k_2$-subspace of the quotient $\F_Q^n/V_1$, and continue.  The subspace-counting theorem gives exactly the displayed product.
\end{proof}

\section{Weighted selection and the Konvalina framework}
Let $w_1,\ldots,w_n$ be commuting weights.  A selection without repetition has weight equal to the product of selected box weights; a selection with repetition is a weakly increasing sequence of boxes with the analogous product weight.

\begin{proposition}[Two universal weighted coefficients]
The total weights are
\begin{align}
C_k(w_1,\ldots,w_n)&=e_k(w_1,\ldots,w_n),\\
S_k(w_1,\ldots,w_n)&=h_k(w_1,\ldots,w_n),
\end{align}
with generating functions
\[
\sum_{k=0}^n C_k t^k=\prod_{j=1}^n(1+w_jt),
\qquad
\sum_{k\ge0}S_k t^k=\prod_{j=1}^n(1-w_jt)^{-1}.
\]
For geometric weights,
\[
C_k(1,q,\ldots,q^{n-1})=q^{\binom k2}\qbinom nk,
\qquad
S_k(1,q,\ldots,q^{n-k})=\qbinom nk.
\]
\end{proposition}
\begin{proof}
Expanding the first product chooses either $1$ or $w_jt$ from every factor, so the coefficient is the sum over $k$-subsets.  Expanding each reciprocal as a geometric series counts multiplicities, hence weak selections.  The geometric specializations are \eqref{eq:elementary} and \eqref{eq:complete}.
\end{proof}

This framework simultaneously explains ordinary binomial coefficients (all weights $1$), Gaussian coefficients (geometric weights), and many Stirling-type triangles (appropriate arithmetic or geometric weight systems).  Their recurrences are simply the alternatives of omitting the last box or using it.

\section{q-Vandermonde convolution and multinomial refinements}
\begin{theorem}[$q$-Vandermonde convolution]
For nonnegative $m,n,r$,
\begin{equation}
\genfrac{[}{]}{0pt}{}{m+n}{r}_q
=\sum_k q^{(m-k)(r-k)}
\genfrac{[}{]}{0pt}{}{m}{k}_q
\genfrac{[}{]}{0pt}{}{n}{r-k}_q,\label{eq:qvand}
\end{equation}
where out-of-range terms vanish.
\end{theorem}
\begin{proof}
Factor the product in the finite $q$-binomial theorem into its first $m$ and last $n$ factors:
\[
\prod_{j=0}^{m+n-1}(1+tq^j)
=\prod_{j=0}^{m-1}(1+tq^j)
 \prod_{j=0}^{n-1}(1+tq^{m+j}).
\]
The coefficient of $t^r$ on the left is $q^{\binom r2}\genfrac{[}{]}{0pt}{}{m+n}{r}_q$.  Choosing degree $k$ from the first product and $r-k$ from the second gives
\[
q^{\binom k2}\genfrac{[}{]}{0pt}{}{m}{k}_q
q^{m(r-k)+\binom{r-k}2}\genfrac{[}{]}{0pt}{}{n}{r-k}_q.
\]
After division by $q^{\binom r2}$, the exponent is
$\binom k2+m(r-k)+\binom{r-k}2-\binom r2=(m-k)(r-k)$.
\end{proof}

\begin{corollary}[Flag convolution]
For $a+b+c=n$,
\[
\genfrac{[}{]}{0pt}{}{n}{a,b,c}_q
=\genfrac{[}{]}{0pt}{}{n}{a}_q
 \genfrac{[}{]}{0pt}{}{n-a}{b}_q
=\genfrac{[}{]}{0pt}{}{n}{c}_q
 \genfrac{[}{]}{0pt}{}{n-c}{b}_q.
\]
\end{corollary}
\begin{proof}
Both expressions telescope to $(q;q)_n/((q;q)_a(q;q)_b(q;q)_c)$.
\end{proof}

\section{q-binomial inversion}
\begin{theorem}[$q$-binomial inversion]
For sequences $(a_n)$ and $(b_n)$ in a commutative $\C(q)$-algebra, the relations
\[
b_n=\sum_{k=0}^n\qbinom nk a_k
\]
are equivalent to
\[
a_n=\sum_{k=0}^n(-1)^{n-k}q^{\binom{n-k}{2}}\qbinom nk b_k.
\]
\end{theorem}
\begin{proof}
Substitute the first relation into the proposed inverse.  The coefficient of $a_j$ is
\[
\sum_{k=j}^n(-1)^{n-k}q^{\binom{n-k}{2}}
\qbinom nk\qbinom kj.
\]
The factorial definition gives
\[
\qbinom nk\qbinom kj=\qbinom nj\genfrac{[}{]}{0pt}{}{n-j}{k-j}_q.
\]
Put $s=n-k$ and $N=n-j$.  The coefficient becomes
\[
\qbinom nj\sum_{s=0}^{N}(-1)^sq^{\binom s2}\genfrac{[}{]}{0pt}{}{N}{s}_q.
\]
By the finite $q$-binomial theorem this sum is $(1;q)_N$, equal to $0$ for $N>0$ and $1$ for $N=0$.  Thus only $j=n$ survives.  Since the two triangular matrices have diagonal entries $1$, either relation implies the other.
\end{proof}

\section{Euler's two series}
\begin{theorem}[Euler expansions]
For $|q|<1$,
\begin{align}
(z;q)_\infty&=\sum_{n=0}^{\infty}rac{(-1)^nq^{\binom n2}}{(q;q)_n}z^n,\label{eq:euler1}\\
\frac1{(z;q)_\infty}&=\sum_{n=0}^{\infty}\frac{z^n}{(q;q)_n},\qquad |z|<1.\label{eq:euler2}
\end{align}
\end{theorem}
\begin{proof}
For \eqref{eq:euler1}, apply the finite theorem and let $N\to\infty$:
\[
(z;q)_N=\sum_{n=0}^N\frac{(-1)^nq^{\binom n2}}{(q;q)_n}
\frac{(q;q)_N}{(q;q)_{N-n}}z^n.
\]
For fixed $n$, the last quotient tends to $1$.  The factor $|q|^{\binom n2}$ supplies uniform summable domination on compact $z$-sets, so dominated convergence is valid.

For \eqref{eq:euler2}, write $F(z)=1/(z;q)_\infty$.  The product equation gives $(1-z)F(z)=F(qz)$.  If $F(z)=\sum c_nz^n$, comparison of coefficients yields $(1-q^n)c_n=c_{n-1}$, with $c_0=1$.  Therefore $c_n=1/(q;q)_n$.  The resulting series converges for $|z|<1$ and has the same functional equation and initial value as the product.
\end{proof}

\section{The infinite q-binomial theorem}
\begin{theorem}[Cauchy--Heine $q$-binomial theorem]
For $|q|<1$ and $|z|<1$,
\begin{equation}
\sum_{n=0}^{\infty}\frac{(a;q)_n}{(q;q)_n}z^n
=\frac{(az;q)_\infty}{(z;q)_\infty}.\label{eq:iqbt}
\end{equation}
The identity also holds formally in $\C(a,q)[[z]]$.
\end{theorem}
\begin{proof}
Let $F(z)=(az;q)_\infty/(z;q)_\infty$.  Cancelling the first factors gives
\[
(1-z)F(z)=(1-az)F(qz).
\]
Writing $F(z)=\sum_{n\ge0}c_nz^n$, with $c_0=1$, and comparing the coefficient of $z^n$ for $n\ge1$ gives
\[
(1-q^n)c_n=(1-aq^{n-1})c_{n-1}.
\]
Thus $c_n=(a;q)_n/(q;q)_n$.  This coefficient recursion proves the formal identity.  Analytically, both sides converge normally on compact subsets of $|z|<1$, so the same computation proves \eqref{eq:iqbt}.
\end{proof}

\begin{corollary}[Finite reciprocal product]
For $m\ge0$,
\[
\frac1{(z;q)_{m+1}}=\sum_{r=0}^{\infty}
\genfrac{[}{]}{0pt}{}{m+r}{r}_q z^r.
\]
Consequently,
\[
\sum_{n=k}^{\infty}\qbinom nk z^n=\frac{z^k}{(z;q)_{k+1}}.
\]
\end{corollary}
\begin{proof}
Set $a=q^{m+1}$ in \eqref{eq:iqbt}; the product ratio telescopes to $1/(z;q)_{m+1}$ and $(q^{m+1};q)_r/(q;q)_r=\genfrac{[}{]}{0pt}{}{m+r}{r}_q$.  Then put $m=k$ and replace $r$ by $n-k$.
\end{proof}

\section{Basic hypergeometric series}
\begin{definition}
The basic hypergeometric series is
\[
{}_r\phi_s\!\left(\begin{matrix}a_1,\ldots,a_r\\ b_1,\ldots,b_s\end{matrix};q,z\right)
:=\sum_{n=0}^{\infty}
\frac{(a_1,\ldots,a_r;q)_n}{(q,b_1,\ldots,b_s;q)_n}
\left((-1)^nq^{\binom n2}\right)^{1+s-r}z^n.
\]
When $r=s+1$, the extra factor has exponent zero.  The $q$-binomial theorem is the evaluation
\[
{}_1\phi_0(a;-;q,z)=\frac{(az;q)_\infty}{(z;q)_\infty}.
\]
\end{definition}

\begin{proposition}[Term ratio and convergence in the balanced case]
For ${}_{s+1}\phi_s$, the ratio of consecutive terms tends to $z$ when no parameter forces termination.  Hence the radius of convergence in $z$ is $1$.
\end{proposition}
\begin{proof}
The ratio is
\[
z\frac{\prod_{i=1}^{s+1}(1-a_iq^n)}{(1-q^{n+1})\prod_{j=1}^{s}(1-b_jq^n)},
\]
which tends to $z$ because $q^n\to0$.  The ratio test proves the assertion.
\end{proof}

\section{Heine's transformation and q-Gauss summation}
\begin{theorem}[Heine transformation]
For $|q|<1$ and parameters in a common domain of absolute convergence,
\begin{equation}
{}_2\phi_1\!\left(\begin{matrix}a,b\\c\end{matrix};q,z\right)
=\frac{(b,az;q)_\infty}{(c,z;q)_\infty}
{}_2\phi_1\!\left(\begin{matrix}c/b,z\\az\end{matrix};q,b\right).
\label{eq:heine}
\end{equation}
Elsewhere it holds by meromorphic continuation.
\end{theorem}
\begin{proof}
Use tail quotients to write
\[
\frac{(b;q)_n}{(c;q)_n}
=\frac{(b;q)_\infty}{(c;q)_\infty}
\frac{(cq^n;q)_\infty}{(bq^n;q)_\infty}.
\]
Apply the $q$-binomial theorem to the last ratio:
\[
\frac{(cq^n;q)_\infty}{(bq^n;q)_\infty}
=\sum_{m\ge0}\frac{(c/b;q)_m}{(q;q)_m}b^mq^{mn}.
\]
Substitute this into the defining ${}_2\phi_1$ and interchange the absolutely convergent sums.  The inner $n$-sum is
\[
\sum_{n\ge0}\frac{(a;q)_n}{(q;q)_n}(zq^m)^n
=\frac{(azq^m;q)_\infty}{(zq^m;q)_\infty}
=\frac{(az;q)_\infty}{(z;q)_\infty}
\frac{(z;q)_m}{(az;q)_m}.
\]
Collecting factors leaves exactly the right side of \eqref{eq:heine}.
\end{proof}

\begin{theorem}[$q$-Gauss summation]
If $|c/(ab)|<1$, then
\begin{equation}
{}_2\phi_1\!\left(\begin{matrix}a,b\\c\end{matrix};q,\frac{c}{ab}\right)
=\frac{(c/a,c/b;q)_\infty}{(c,c/(ab);q)_\infty}.
\label{eq:qgauss}
\end{equation}
\end{theorem}
\begin{proof}
Set $z=c/(ab)$ in Heine's transformation.  Since $az=c/b$, one upper parameter of the transformed series cancels its lower parameter, and the remaining series is
\[
{}_1\phi_0\!\left(c/(ab);-;q,b\right)
=\frac{(c/a;q)_\infty}{(b;q)_\infty}
\]
by the $q$-binomial theorem.  Multiplication by Heine's prefactor $(b,c/b;q)_\infty/(c,c/(ab);q)_\infty$ gives \eqref{eq:qgauss}.
\end{proof}

\section{Terminating q-Chu--Vandermonde summations}
\begin{theorem}[Two forms]
For $n\in\N$,
\begin{align}
{}_2\phi_1\!\left(\begin{matrix}q^{-n},a\\c\end{matrix};q,\frac{cq^n}{a}\right)
&=\frac{(c/a;q)_n}{(c;q)_n},\label{eq:qchu1}\\
{}_2\phi_1\!\left(\begin{matrix}q^{-n},a\\c\end{matrix};q,q\right)
&=a^n\frac{(c/a;q)_n}{(c;q)_n}.\label{eq:qchu2}
\end{align}
\end{theorem}
\begin{proof}
For \eqref{eq:qchu1}, put $b=q^{-n}$ in $q$-Gauss.  The left side terminates at $n$, while the right side simplifies by tail cancellation:
\[
\frac{(c/a,cq^n;q)_\infty}{(c,cq^n/a;q)_\infty}
=\frac{(c/a;q)_n}{(c;q)_n}.
\]
Initially this derivation is valid on a nonempty open set satisfying the convergence restrictions after approaching $q^{-n}$ by generic $b$.  Both sides of the resulting terminated identity are rational functions of $a$ and $c$, so equality extends identically.

To derive \eqref{eq:qchu2}, reverse the finite sum in \eqref{eq:qchu1}.  The elementary reversal formulas
\[
(q^{-n};q)_{n-j}=\frac{(q^{-n};q)_n}{(q^{-j};q)_j},
\qquad
(x;q)_{n-j}=\frac{(x;q)_n}{(xq^{n-j};q)_j},
\]
together with \eqref{eq:reverse}, reduce the reversed identity, after cancelling the common factor $(a;q)_n/(c;q)_n$, precisely to \eqref{eq:qchu2}.  Every manipulation is between finite products, so no convergence issue remains.
\end{proof}

\section{The Jacobi triple product}
\begin{lemma}[Finite triple-product identity]
For $N\ge0$,
\begin{equation}
(z;q)_N(q/z;q)_N
=\sum_{k=-N}^{N}(-1)^kq^{\binom k2}z^k
\genfrac{[}{]}{0pt}{}{2N}{N+k}_q.
\label{eq:finitejtp}
\end{equation}
\end{lemma}
\begin{proof}
Apply the finite $q$-binomial theorem to both products and compare the coefficient of $z^k$ in their product.  For fixed $k$, the resulting convolution is a $q$-Vandermonde sum; inserting \eqref{eq:qvand} with the corresponding shifted indices reduces it to $\genfrac{[}{]}{0pt}{}{2N}{N+k}_q$, with the displayed sign and power.  This proves every Laurent coefficient and hence the identity.
\end{proof}

\begin{theorem}[Jacobi triple product]
For $|q|<1$ and $z\ne0$,
\begin{equation}
(q;q)_\infty(z;q)_\infty(q/z;q)_\infty
=\sum_{k\in\Z}(-1)^kq^{\binom k2}z^k.
\label{eq:jtp}
\end{equation}
The convergence is normal on compact annuli in the $z$-plane.
\end{theorem}
\begin{proof}
For fixed $k$,
\[
\genfrac{[}{]}{0pt}{}{2N}{N+k}_q
=\frac{(q;q)_{2N}}{(q;q)_{N+k}(q;q)_{N-k}}
\longrightarrow\frac1{(q;q)_\infty}.
\]
Multiply \eqref{eq:finitejtp} by $(q;q)_\infty$ and let $N\to\infty$.  The product side converges locally uniformly.  On any compact annulus $r\le|z|\le R$, the bilateral summand is bounded by a constant times $|q|^{k^2/3}(R^{|k|}+r^{-|k|})$ for all sufficiently large $|k|$, which is summable.  Dominated convergence therefore gives \eqref{eq:jtp}.
\end{proof}

\begin{corollary}[Euler pentagonal theorem]
\[
(q;q)_\infty=\sum_{k\in\Z}(-1)^kq^{k(3k-1)/2}.
\]
\end{corollary}
\begin{proof}
In \eqref{eq:jtp}, replace $q$ by $q^3$ and set $z=q$.  The product on the left becomes
$(q^3;q^3)_\infty(q;q^3)_\infty(q^2;q^3)_\infty=(q;q)_\infty$, while the exponent on the right is $3\binom k2+k=k(3k-1)/2$.
\end{proof}

\section{Roots of unity and the q-Lucas theorem}
\begin{theorem}[$q$-Lucas congruence]
Let $d\ge1$, and write $n=ad+b$, $k=rd+s$ with $0\le b,s<d$.  Then
\[
\genfrac{[}{]}{0pt}{}{n}{k}_q
\equiv \binom ar\genfrac{[}{]}{0pt}{}{b}{s}_q
\pmod{\Phi_d(q)}.
\]
Equivalently, at every primitive $d$th root $\zeta$,
\[
\genfrac{[}{]}{0pt}{}{n}{k}_{\zeta}
=\binom ar\genfrac{[}{]}{0pt}{}{b}{s}_{\zeta}.
\]
\end{theorem}
\begin{proof}
At $q=\zeta$, break the finite product theorem into $a$ full blocks and one partial block:
\[
\prod_{j=0}^{n-1}(1+t\zeta^j)
=\left(\prod_{j=0}^{d-1}(1+t\zeta^j)\right)^a
 \prod_{j=0}^{b-1}(1+t\zeta^j).
\]
Because the numbers $\zeta^j$ are the roots of $x^d-1$,
\[
\prod_{j=0}^{d-1}(1+t\zeta^j)=1-(-t)^d.
\]
Extracting the coefficient of $t^{rd+s}$ and using the finite theorem on the partial block gives
\[
\zeta^{\binom{k}{2}}\genfrac{[}{]}{0pt}{}{n}{k}_{\zeta}
=(-1)^{r(d+1)}\binom ar\,
\zeta^{\binom{s}{2}}\genfrac{[}{]}{0pt}{}{b}{s}_{\zeta}.
\]
Now
\[
\zeta^{\binom{k}{2}-\binom{s}{2}}
=(-1)^{r(rd+2s-1)}=(-1)^{r(d+1)},
\]
because the difference between the two exponents of $-1$ is $r(d(r-1)+2s-2)$, an even integer.  Cancelling the common phase proves the evaluation.  A polynomial vanishes at all primitive $d$th roots exactly when it is divisible by $\Phi_d$, proving the congruence.
\end{proof}

\section{Cyclic sieving for rotating subsets}
\begin{theorem}
Let $C_n$ act on the $k$-subsets of $\{0,1,\ldots,n-1\}$ by cyclic rotation.  Let $g$ be rotation by one step and put $X(q)=\qbinom nk$.  For every integer $m$,
\[
X(e^{2\pi i m/n})=\#\{S:g^mS=S\}.
\]
Thus $(X,C_n,X(q))$ exhibits the cyclic sieving phenomenon.
\end{theorem}
\begin{proof}
The permutation $g^m$ has $G=\gcd(n,m)$ cycles, each of length $L=n/G$.  A fixed subset must be a union of whole cycles.  Hence the number of fixed $k$-subsets is $\binom{G}{k/L}$ if $L\mid k$, and zero otherwise.

The number $\zeta=e^{2\pi im/n}$ has order $L$.  Write $n=GL$ and $k=aL+b$, $0\le b<L$.  The $q$-Lucas theorem gives
\[
\qbinom nk\bigg|_{q=\zeta}
=\binom Ga\genfrac{[}{]}{0pt}{}{0}{b}_{\zeta}.
\]
This is $\binom Ga$ when $b=0$ and $0$ otherwise, exactly the fixed-point count.
\end{proof}

\section{The quantum-plane binomial theorem}
\begin{theorem}
Let $x,y$ belong to an associative algebra and satisfy $yx=qxy$, with central $q$.  Then
\[
(x+y)^n=\sum_{k=0}^n\qbinom nk x^{n-k}y^k.
\]
\end{theorem}
\begin{proof}
Induct on $n$, multiplying on the right by $x+y$.  Since $y^kx=q^kxy^k$, the coefficient of $x^{n+1-k}y^k$ becomes
\[
q^k\qbinom nk+\qbinom n{k-1}=\genfrac{[}{]}{0pt}{}{n+1}{k}_q
\]
by \eqref{eq:pascal2}.  The initial case is immediate.
\end{proof}

This is not merely an analogy: Gaussian coefficients are exactly the structure constants required to expand powers on the quantum plane.  Iterating the proof gives a $q$-multinomial theorem for variables satisfying $x_jx_i=qx_ix_j$ whenever $i<j$.

\section{q-differentiation, q-exponentials, and Jackson integration}
Assume $q\ne1$.
\begin{definition}
The Jackson $q$-derivative is
\[
D_qf(x):=\frac{f(x)-f(qx)}{(1-q)x},
\]
with the value at $x=0$ supplied by continuity when it exists.
\end{definition}

\begin{proposition}
For $n\ge0$,
\[
D_qx^n=[n]_qx^{n-1},\qquad
D_q^kx^n=\frac{[n]_q!}{[n-k]_q!}x^{n-k}.
\]
Moreover
\[
D_q(fg)(x)=g(x)D_qf(x)+f(qx)D_qg(x).
\]
\end{proposition}
\begin{proof}
The monomial formula follows from $(x^n-q^nx^n)/((1-q)x)$.  Iteration gives the second formula.  For the product rule, insert and subtract $f(qx)g(x)$ in $f(x)g(x)-f(qx)g(qx)$ and divide by $(1-q)x$.
\end{proof}

\begin{definition}[Two $q$-exponentials]
\begin{align*}
e_q(z)&:=\sum_{n\ge0}\frac{z^n}{[n]_q!}
      =\frac1{((1-q)z;q)_\infty},\\
E_q(z)&:=\sum_{n\ge0}\frac{q^{\binom n2}z^n}{[n]_q!}
      =(-(1-q)z;q)_\infty.
\end{align*}
\end{definition}
\begin{proof}[Proof of the product representations]
Since $[n]_q!=(q;q)_n/(1-q)^n$, the two series are Euler's reciprocal and direct expansions with $z$ replaced by $(1-q)z$ and $-(1-q)z$, respectively.
\end{proof}

\begin{corollary}
\[
e_q(z)E_q(-z)=1,
\qquad D_qe_q(\lambda x)=\lambda e_q(\lambda x).
\]
\end{corollary}
\begin{proof}
The product identity follows by cancelling identical infinite products.  Termwise differentiation and $D_qx^n=[n]_qx^{n-1}$ prove the eigenfunction equation inside the disk of convergence, and then by analytic continuation wherever appropriate.
\end{proof}

\begin{definition}[Jackson integral]
For $0<q<1$ and suitable $f$,
\[
\int_0^a f(x)\,d_qx:=(1-q)a\sum_{j=0}^{\infty}q^jf(aq^j).
\]
\end{definition}

\begin{theorem}[Fundamental theorem of $q$-calculus]
If $F$ is continuous at $0$ and the series converges, then
\[
\int_0^a D_qF(x)\,d_qx=F(a)-F(0).
\]
\end{theorem}
\begin{proof}
Substitution gives a telescoping series:
\[
(1-q)a\sum_{j\ge0}q^j
\frac{F(aq^j)-F(aq^{j+1})}{(1-q)aq^j}
=\sum_{j\ge0}\bigl(F(aq^j)-F(aq^{j+1})\bigr).
\]
Its $N$th partial sum is $F(a)-F(aq^{N+1})$, which tends to $F(a)-F(0)$.
\end{proof}

\section{The q-gamma function and classical limits}
For $0<q<1$ and $x$ away from the poles, define
\[
\Gamma_q(x):=(1-q)^{1-x}\frac{(q;q)_\infty}{(q^x;q)_\infty}.
\]
\begin{proposition}[Functional equation]
\[
\Gamma_q(x+1)=[x]_q\Gamma_q(x),\qquad \Gamma_q(1)=1.
\]
\end{proposition}
\begin{proof}
Use $(q^x;q)_\infty=(1-q^x)(q^{x+1};q)_\infty$ and simplify:
\[
\frac{\Gamma_q(x+1)}{\Gamma_q(x)}
=(1-q)^{-1}(1-q^x)=[x]_q.
\]
The value at $1$ is immediate.
\end{proof}

\begin{proposition}[Classical finite limits]
For fixed $n,k$ and complex $x$,
\[
\lim_{q\to1}[n]_q=n,
\qquad
\lim_{q\to1}\qbinom nk=\binom nk,
\qquad
\lim_{q\to1}\frac{(q^x;q)_n}{(1-q)^n}=(x)_n.
\]
\end{proposition}
\begin{proof}
For every fixed $j$, $(1-q^{x+j})/(1-q)\to x+j$ by differentiation at $q=1$ or the exponential parametrization $q=e^{-t}$.  Multiplying finitely many factors proves all three limits.
\end{proof}

\begin{corollary}[Classical hypergeometric limit]
For fixed parameters and $|z|<1$,
\[
{}_{r}\phi_{r-1}\!\left(\begin{matrix}q^{a_1},\ldots,q^{a_r}\\q^{b_1},\ldots,q^{b_{r-1}}\end{matrix};q,z\right)
\longrightarrow
{}_{r}F_{r-1}\!\left(\begin{matrix}a_1,\ldots,a_r\\b_1,\ldots,b_{r-1}\end{matrix};z\right).
\]
\end{corollary}
\begin{proof}
The preceding finite-factor limit gives termwise convergence.  On compact subdisks of $|z|<1$, term ratios are uniformly bounded by a number below $1$ for $q$ sufficiently close to $1$, so dominated convergence interchanges limit and sum.
\end{proof}

\section{Asymptotic expansion near q=1}
Put $q=e^{-t}$ with $t\downarrow0$.  The elementary expansion
\[
\log\frac{1-e^{-x}}x=-\frac x2+
\sum_{r=1}^{\infty}\frac{B_{2r}}{2r(2r)!}x^{2r}
\]
is convergent for $|x|<2\pi$ and is also an asymptotic expansion.

\begin{theorem}[All-orders fixed-parameter expansion]
For fixed $0\le k\le n$,
\begin{align}
\log\frac{\genfrac{[}{]}{0pt}{}{n}{k}_{e^{-t}}}{\binom nk}
={}&-\frac{k(n-k)}2t\\
&+\sum_{r=1}^{\infty}\frac{B_{2r}t^{2r}}{2r(2r)!}
\left(S_{2r}(n)-S_{2r}(k)-S_{2r}(n-k)\right),\label{eq:allorder}
\end{align}
where $S_m(N)=\sum_{j=1}^{N}j^m$.  In particular,
\[
\genfrac{[}{]}{0pt}{}{n}{k}_{e^{-t}}
=\binom nk\exp\left(
-\frac{k(n-k)}2t+
\frac{k(n-k)(n+1)}{24}t^2+O(t^4)\right).
\]
\end{theorem}
\begin{proof}
Write the Gaussian coefficient as
\[
\prod_{j=n-k+1}^{n}\frac{1-e^{-jt}}{1-e^{-(j-n+k)t}},
\]
or, more transparently, take the logarithm of the quotient of $[n]_q!$ by $[k]_q![n-k]_q!$.  Apply the displayed Bernoulli expansion to each fixed factor and subtract the three power sums.  The linear power-sum difference is $k(n-k)$; the quadratic difference is
$S_2(n)-S_2(k)-S_2(n-k)=k(n-k)(n+1)$.  This gives \eqref{eq:allorder} and its truncation.
\end{proof}

\begin{theorem}[Dilogarithmic expansion of an infinite product]
For fixed $|a|<1$,
\[
\log(a;e^{-t})_\infty
=-\frac{\Li_2(a)}t+\frac12\log(1-a)
-\frac{t}{12}\frac{a}{1-a}+O(t^3)
\quad(t\downarrow0).
\]
More generally, an all-orders expansion follows by inserting the Bernoulli expansion of $(1-e^{-tr})^{-1}$ into the logarithmic series.
\end{theorem}
\begin{proof}
From the logarithmic expansion,
\[
\log(a;e^{-t})_\infty=-\sum_{r\ge1}\frac{a^r}{r(1-e^{-tr})}.
\]
Uniformly for $|a|\le\rho<1$,
\[
\frac1{1-e^{-tr}}=\frac1{tr}+\frac12+\frac{tr}{12}+O(t^3r^3).
\]
After multiplication by $a^r/r$, the error is summable.  The three sums are respectively $\Li_2(a)$, $-\log(1-a)$, and $a/(1-a)$, with the displayed signs.
\end{proof}

\section{Fixed-q growth and central coefficients}
\begin{proposition}
For fixed $|q|<1$ and fixed integer $r$,
\[
\lim_{n\to\infty}\genfrac{[}{]}{0pt}{}{2n}{n+r}_q=\frac1{(q;q)_\infty}.
\]
For real $q>1$,
\[
\genfrac{[}{]}{0pt}{}{2n}{n}_q
=q^{n^2}\genfrac{[}{]}{0pt}{}{2n}{n}_{q^{-1}}
\sim \frac{q^{n^2}}{(q^{-1};q^{-1})_\infty}.
\]
\end{proposition}
\begin{proof}
For $|q|<1$, use the factorial quotient:
\[
\genfrac{[}{]}{0pt}{}{2n}{n+r}_q
=\frac{(q;q)_{2n}}{(q;q)_{n+r}(q;q)_{n-r}},
\]
and take limits of all three products.  For $q>1$, reciprocity contributes the factor $q^{(n)(n)}$, after which the first result applies to $q^{-1}$.
\end{proof}

\section{Alternating sums and interpolation at a q-grid}
The finite theorem gives immediately
\begin{equation}
\sum_{k=0}^{n}(-1)^kq^{\binom k2}\qbinom nk x^k=(x;q)_n.
\label{eq:altsum}
\end{equation}
This compact identity contains a family of orthogonality relations.

\begin{corollary}[Exponential moments]
For $0\le j<n$,
\[
\sum_{k=0}^{n}(-1)^kq^{\binom k2-jk}\qbinom nk=0.
\]
For $j=n$ the same sum equals $(q^{-n};q)_n$.
\end{corollary}
\begin{proof}
Set $x=q^{-j}$ in \eqref{eq:altsum}.  If $j<n$, the product $(q^{-j};q)_n$ contains the factor $1-q^0$; the case $j=n$ is simply the unevaluated product.
\end{proof}

\begin{proposition}[Newton basis on a geometric grid]
The polynomials $(x;q)_0,\ldots,(x;q)_n$ form a basis of the polynomials of degree at most $n$.  Their zeros are nested geometric grids, so coefficients in this basis can be recovered by successive divided differences at $1,q^{-1},q^{-2},\ldots$.
\end{proposition}
\begin{proof}
The polynomial $(x;q)_k$ has degree $k$ and nonzero leading coefficient.  The change-of-basis matrix from these polynomials to $1,x,\ldots,x^n$ is triangular with nonzero diagonal, hence invertible.  The nested-zero statement follows from the explicit roots; ordinary Newton interpolation on those nodes gives the coefficient recovery.
\end{proof}

\section{Rogers--Szego polynomials}
Define
\[
H_n(z;q):=\sum_{k=0}^n\qbinom nk z^k.
\]
\begin{proposition}[Three-term recurrence]
For $n\ge1$,
\[
H_{n+1}(z;q)=(1+z)H_n(z;q)-z(1-q^n)H_{n-1}(z;q).
\]
\end{proposition}
\begin{proof}
The coefficient of $z^k$ on the right is
\[
\qbinom nk+\qbinom n{k-1}-(1-q^n)\genfrac{[}{]}{0pt}{}{n-1}{k-1}_q.
\]
Use
$\qbinom n{k-1}=\genfrac{[}{]}{0pt}{}{n-1}{k-1}_q+q^{n-k+1}\genfrac{[}{]}{0pt}{}{n-1}{k-2}_q$
and the corresponding Pascal relations, or reduce all terms to a common factorial denominator; the result is $\genfrac{[}{]}{0pt}{}{n+1}{k}_q$.
\end{proof}

\section{Principal specialization of symmetric functions}
The identities \eqref{eq:elementary}--\eqref{eq:complete} are the one-row and one-column cases of a broad principle: principal specialization converts symmetric functions into product formulas built from $q$-Pochhammer symbols.

\begin{theorem}[Principal specialization of a Schur polynomial]
Let $\lambda=(\lambda_1,\ldots,\lambda_m)$ be a partition padded with zeros, and put $n(\lambda)=\sum_i(i-1)\lambda_i$.  Then
\[
s_\lambda(1,q,\ldots,q^{m-1})
=q^{n(\lambda)}\prod_{1\le i<j\le m}
\frac{1-q^{\lambda_i-\lambda_j+j-i}}{1-q^{j-i}}.
\]
\end{theorem}
\begin{proof}
Use the alternant definition
\[
s_\lambda(x_1,\ldots,x_m)=
\frac{\det(x_i^{\lambda_j+m-j})_{i,j}}
     {\det(x_i^{m-j})_{i,j}}.
\]
At $x_i=q^{i-1}$, each determinant is a Vandermonde determinant in the numbers $q^{\lambda_j+m-j}$ or $q^{m-j}$.  Thus the quotient is a product over $i<j$ of differences.  Factor the lower powers of $q$ from each difference.  The total residual power is $n(\lambda)$, and the remaining ratios are exactly those displayed.
\end{proof}

\section{Further connections}
\subsection{MacMahon's q-Catalan quotient}
The polynomial
\[
C_n(q):=\frac1{[n+1]_q}\genfrac{[}{]}{0pt}{}{2n}{n}_q
\]
is a standard $q$-analogue of the Catalan number.  Cyclotomic factorization proves polynomiality: subtracting the exponent of each $\Phi_d$ in $[n+1]_q$ from that in the central Gaussian coefficient gives a nonnegative integer.  At $q=1$ it tends to $\frac1{n+1}\binom{2n}{n}$.

\subsection{Mahonian statistics}
The word model is the two-letter case of the fact that the Gaussian multinomial is the inversion enumerator of words with prescribed multiplicities:
\[
\genfrac{[}{]}{0pt}{}{n}{k_1,\ldots,k_r}_q
=\sum_w q^{\inv(w)}.
\]
The proof repeats the last-letter recurrence, adding the number of letters larger than the appended letter.  For permutations, analogous insertion recurrences lead to the $q$-factorial $[n]_q!$ as the inversion generating polynomial.

\subsection{Partition products}
Euler's reciprocal identity says $1/(q;q)_\infty=\sum_{n\ge0}p(n)q^n$: in expanding $\prod_{j\ge1}(1-q^j)^{-1}$, the exponent chosen from the $j$th factor records the multiplicity of the part $j$.  The direct Euler expansion and the triple product convert product restrictions into signed or bilateral sums.  This is the elementary gateway to the Rogers--Ramanujan identities, affine root-system denominator identities, and modular forms.

\subsection{Formal proof and computer algebra}
All terminating identities in this article live in rational-function fields and can be checked by recurrence, creative telescoping, or normalization to a common product.  This makes them particularly suitable for formalization: finite products and coefficient extraction avoid analytic side conditions, while the analytic theorems can be layered on later through uniform convergence.  Basic hypergeometric summands are $q$-hypergeometric in the algorithmic sense because the quotient of successive terms is rational in $q^n$.

\section{A compact identity atlas}
For reference, the most frequently used formulas are gathered here:
\begin{longtable}{>{\raggedright\arraybackslash}p{0.29\textwidth}p{0.64\textwidth}}
\toprule
Name & Formula\\
\midrule
Concatenation & $(a;q)_{m+n}=(a;q)_m(aq^m;q)_n$\\[1mm]
Reversal & $(a;q)_n=(-a)^nq^{\binom n2}(q^{1-n}/a;q)_n$\\[1mm]
Finite $q$-binomial & $(z;q)_n=\sum_k(-1)^kq^{\binom k2}\qbinom nkz^k$\\[1mm]
Euler I & $(z;q)_\infty=\sum_{n\ge0}(-1)^nq^{\binom n2}z^n/(q;q)_n$\\[1mm]
Euler II & $(z;q)_\infty^{-1}=\sum_{n\ge0}z^n/(q;q)_n$\\[1mm]
Infinite $q$-binomial & $(az;q)_\infty/(z;q)_\infty=\sum_{n\ge0}(a;q)_nz^n/(q;q)_n$\\[1mm]
$q$-Vandermonde & $\genfrac{[}{]}{0pt}{}{m+n}{r}_q=\sum_kq^{(m-k)(r-k)}\genfrac{[}{]}{0pt}{}{m}{k}_q\genfrac{[}{]}{0pt}{}{n}{r-k}_q$\\[1mm]
$q$-binomial inversion & $b_n=\sum_k\qbinom nka_k\iff a_n=\sum_k(-1)^{n-k}q^{\binom{n-k}2}\qbinom nkb_k$\\[1mm]
Heine & Equation \eqref{eq:heine}\\[1mm]
$q$-Gauss & Equation \eqref{eq:qgauss}\\[1mm]
$q$-Lucas & $\genfrac{[}{]}{0pt}{}{ad+b}{rd+s}_q\equiv\binom ar\genfrac{[}{]}{0pt}{}{b}{s}_q\pmod{\Phi_d(q)}$\\[1mm]
Triple product & $(q,z,q/z;q)_\infty=\sum_{k\in\Z}(-1)^kq^{\binom k2}z^k$\\
\bottomrule
\end{longtable}

\section{Proof strategies and diagnostic checks}
A useful hierarchy of methods emerges from the preceding proofs.
\begin{enumerate}[label=\arabic*.]
\item A finite product identity should first be attacked by partitioning exponent sets or by induction using a $q$-Pascal recurrence.
\item A polynomial identity should be interpreted by words, paths, partitions, or subspaces; positivity and reciprocity then become structural rather than accidental.
\item A unilateral infinite series whose coefficient ratio is rational in $q^n$ should be tested against a first-order $q$-difference equation.
\item A ${}_2\phi_1$ transformation often follows by expanding one ratio of tails with the $q$-binomial theorem and interchanging two sums.
\item A root-of-unity evaluation should be organized into complete cyclotomic blocks; $q$-Lucas is the canonical outcome.
\item Limits as $q\to1$ should be taken after extracting powers of $1-q$; all-orders expansions follow from Bernoulli numbers, while moving parameters may introduce dilogarithms.
\end{enumerate}
These methods also provide checks.  Every proposed identity should pass degree, reciprocity, $q=0$, $q=1$, and small-index tests.  For hypergeometric sums, matching the ratio of successive summands and testing terminating parameter values catches most normalization errors.

\section{Historical and bibliographic notes}
The notation $(a;q)_n$ is modern, but the underlying products and series run through work of Euler, Gauss, Cauchy, Heine, and Jacobi.  Gaussian coefficients arose from finite-field counting and from polynomial analogues of ordinary binomial coefficients.  Their interpretations by partitions, inversions, and subspaces are now standard.  The weighted-selection perspective developed by Konvalina places binomial, Stirling, and Gaussian triangles in a common framework.  Modern investigations include cyclotomic and supercongruence phenomena, alternating Gaussian sums, cyclic sieving, formal verification, and specialized central $q$-series identities.  The source archive supplied with this project was used as a guide to those directions; notation and proofs here have been normalized into a single self-contained development.

\begin{thebibliography}{99}
\bibitem{AndrewsPartitions}
G.~E. Andrews, \emph{The Theory of Partitions}, Addison--Wesley, 1976; Cambridge University Press reprint, 1998.

\bibitem{AAR}
G.~E. Andrews, R.~Askey, and R.~Roy, \emph{Special Functions}, Cambridge University Press, 1999.

\bibitem{Bressoud}
D.~M. Bressoud, \emph{Proofs and Confirmations: The Story of the Alternating Sign Matrix Conjecture}, Cambridge University Press, 1999.

\bibitem{DLMF}
NIST Digital Library of Mathematical Functions, Chapter 17, ``$q$-Hypergeometric and Related Functions,'' current online edition.

\bibitem{Fine}
N.~J. Fine, \emph{Basic Hypergeometric Series and Applications}, American Mathematical Society, 1988.

\bibitem{GasperRahman}
G.~Gasper and M.~Rahman, \emph{Basic Hypergeometric Series}, 2nd ed., Cambridge University Press, 2004.

\bibitem{KacCheung}
V.~Kac and P.~Cheung, \emph{Quantum Calculus}, Springer, 2002.

\bibitem{Konvalina}
J.~Konvalina, ``A unified interpretation of the binomial coefficients, the Stirling numbers, and the Gaussian coefficients,'' \emph{American Mathematical Monthly} 107 (2000), 901--910.

\bibitem{Kupershmidt}
B.~A. Kupershmidt, ``$q$-Newton binomial: from Euler to Gauss,'' \emph{Journal of Nonlinear Mathematical Physics} 7 (2000), 244--262.

\bibitem{Macdonald}
I.~G. Macdonald, \emph{Symmetric Functions and Hall Polynomials}, 2nd ed., Oxford University Press, 1995.

\bibitem{RSW}
V.~Reiner, D.~Stanton, and D.~White, ``The cyclic sieving phenomenon,'' \emph{Journal of Combinatorial Theory, Series A} 108 (2004), 17--50.

\bibitem{Stanley1}
R.~P. Stanley, \emph{Enumerative Combinatorics}, Vol.~1, 2nd ed., Cambridge University Press, 2012.

\bibitem{Stanley2}
R.~P. Stanley, \emph{Enumerative Combinatorics}, Vol.~2, 2nd ed., Cambridge University Press, 2023.

\bibitem{Zeilberger}
M.~Petkovsek, H.~S. Wilf, and D.~Zeilberger, \emph{A=B}, A K Peters, 1996.
\end{thebibliography}
\end{document}
